\documentclass[DM,lsstdraft,authoryear,toc]{lsstdoc}
\input{meta}

% Package imports go here.

% Local commands go here.

%If you want glossaries
%\input{aglossary.tex}
%\makeglossaries

\title{Implementation of Incremental Template Generation}

% This can write metadata into the PDF.
% Update keywords and author information as necessary.
\hypersetup{
    pdftitle={Implementation of Incremental Template Generation},
    pdfauthor={Eric C. Bellm, Erin L. Howard},
    pdfkeywords={}
}

% Optional subtitle
% \setDocSubtitle{A subtitle}

\input{authors}

\setDocRef{DMTN-344}
\setDocUpstreamLocation{\url{https://github.com/lsst-dm/dmtn-344}}

\date{\vcsDate}

% Optional: name of the document's curator
% \setDocCurator{The Curator of this Document}

\setDocAbstract{%
Describes the process of incremental template generation used by Alert Production prior to DR1.
}

% Change history defined here.
% Order: oldest first.
% Fields: VERSION, DATE, DESCRIPTION, OWNER NAME.
% See LPM-51 for version number policy.
\setDocChangeRecord{%
  \addtohist{1}{YYYY-MM-DD}{Unreleased.}{Bellm}
}


\begin{document}

% Create the title page.
\maketitle
% Frequently for a technote we do not want a title page  uncomment this to remove the title page and changelog.
% use \mkshorttitle to remove the extra pages

% ADD CONTENT HERE
% You can also use the \input command to include several content files.

\section{Motivation for Early Template Generation} \label{sec:motivation}

\citeds{DMTN-107}

\section{Constraints and Boundary Conditions} \label{sec: constraints}

tract/patch inputs

APDB (lightcurves, not just diffim)

slow vetting and deployment

disk space

avoiding baseline changes

\section{Approaches to Early Templates} \label{sec:approaches}

\subsection{Data Release Processing} \label{sec:drp}

During commissioning, small biweekly Data Release Processing campaigns run on the available data were sufficient to generate early templates for Deep Drilling Fields (DDFs) as well as specific fields that received frequent pointed observations.
The initial DDF templates deployed were generated by such campaigns.
Likewise, we deployed about 1k square degrees of templates processed as part of Data Preview 2.

However, as the rate of data taking increased, these test campaigns could no longer operate over all available data.
Accordingly, we developed specialized processing pipelines specifically for incremental template generation (\S \ref{sec:incremental}).

In steady state operations, annual data releases will be the source of updated templates, beginning with Data Release 1.

\subsection{Incremental Rubin Templates} \label{sec:incremental}

Incremental template processing uses a subset of the Data Release pipelines.
In particular, it omits the global photometric and astrometric calibration steps that require large amount of well-dithered input data as well as many of the later measurement stages.

\subsection{DECam Templates} \label{sec:decam}

DECam coadds are available over a large fraction of the Southern Hemisphere sky.
An adaptor for using such coadds as templates in the Rubin pipelines is available, and has been used for initial processing of some Target-of-Opportunity data.
However, head-to-head comparisons between subtractions of the same data against Rubin and DECam templates shows that the latter can exhibit order-of-magnitude higher false-positive rates, too high for production deployment against the APDB.



\section{Challenges Encountered} \label{sec:challenges}

\begin{enumerate}
\item varying image quality
\item irregular data-taking in early ops
\item inconsistent dithering
\item scaling up vetting
\end{enumerate}


\section{Implementing Incremental Template Generation} \label{sec:implementation}

\subsection{Input Selection}

\subsection{Vetting Metrics}

nimages
gaps
PSF variations

\section{Forecasting Year 1 Templates} \label{sec:forecast}



%\appendix

\section{Acknowledgements}

This material is based upon work supported in part by the National Science Foundation through Cooperative Agreements AST-1258333 and AST-2241526 and Cooperative Support Agreements AST-1202910 and AST-2211468 managed by the Association of Universities for Research in Astronomy (AURA), and the Department of Energy under Contract No.\ DE-AC02-76SF00515 with the SLAC National Accelerator Laboratory managed by Stanford University.
Additional Rubin Observatory funding comes from private donations, grants to universities, and in-kind support from LSST-DA Institutional Members.

% Include all the relevant bib files.
% https://lsst-texmf.lsst.io/lsstdoc.html#bibliographies
\section{References} \label{sec:bib}
\renewcommand{\refname}{} % Suppress default Bibliography section
\bibliography{local,lsst,lsst-dm,refs_ads,refs,books}

% Make sure lsst-texmf/bin/generateAcronyms.py is in your path
%\section{Acronyms} \label{sec:acronyms}
%\input{acronyms.tex}
% If you want glossary uncomment below -- comment out the two lines above
%\printglossaries





\end{document}
